\documentclass[11pt]{article}
\usepackage[T1]{fontenc}
\usepackage{lmodern,amsmath,amssymb,amsthm}
\usepackage[margin=1in]{geometry}
\newtheorem{proposition}{Proposition}
\newcommand{\E}{\mathbb E}
\newcommand{\Z}{\mathbb Z}
\title{Difference decomposition and the additive-quadruple estimate}
\author{/root/sol\_protocol\_final\_review}
\date{7 November 2026}
\begin{document}
\maketitle

\section{Inputs and conventions}
We use the natural-ultrapower convention of GTZ Appendix A and the
project shift theorem in tautification-003. Its acceptance remains
separate. A bound $\gg1$ is positive and standard. Write $N$ for an
unbounded limit integer, $J=[N]$ and $I=[-N,N]\cap{}^*\Z$. All changes
of variables are integer identities at the representative indices.
The author erratum guides the proof. Its Proposition 3.1 and modified
Proposition 7.3 are the conclusions here, rather than assumed premises.

The shift theorem gives a lower-support sequence with fixed geometry
and uniform Lipschitz bound. It gives neither a separated finite
decomposition nor a correlation conclusion. Products, conjugates,
and translates use product manifolds, conjugated tests, and translated
polynomial orbits, respectively. A bounded internal scalar depending
only on parameters is degree zero in the averaging variable.

\section{Difference decomposition}
\begin{proposition}[LP02]
For standard $s\ge3$, let $\chi$ be a multidegree $(1,s-1)$
nilcharacter. For every standard $\epsilon>0$ each component of
$D_k(h,n)=\chi(h+k,n)\otimes\overline{\chi(h,n)}$
is approximated uniformly in $k,h,n$ to error $\epsilon$ by
\[
 \sum_{j=1}^{L} \Psi_{j,k,a_j}(n) B_{j,k}(h,n).
\]
Here $L$ is standard, $\Psi_{j,k,a_j}$ is a component of a
unit-vector degree-$(s-1)$ nilcharacter, and $B_{j,k}$ is supported
on the downset below $(1,s-2)$. The families are internal with
standard finite geometry lists and uniform bounds.
\end{proposition}
\begin{proof}
Apply the project shift theorem, Theorem 3, with positive index
$(1,0)$ and shift $(k,0)$. The support is in the union of the downsets
$J_0=\downarrow(0,s-1)$ and $J_1=\downarrow(1,s-2)$.
We now give the fixed geometric split explicitly. Write $G/\Gamma$
for the taut lower-support representation from the shift theorem and
$J=J_0\cup J_1$. For each positive $a\in J$ take a rational basis of
$\mathfrak g_a$ for the rational structure of $\Gamma$. This exists
because $G_a$ is rational. Take a free rational Lie algebra with one
generator of weight $a$ for each basis vector. Kill all homogeneous
words whose total weight is outside $J$. The excluded weights form
an upward-closed set, so these words generate a homogeneous ideal.
All generator weights are positive. Thus word length is bounded by
$\max_{a\in J}|a|=s-1$ in the quotient, which is finite-dimensional
and nilpotent. Let $\mathfrak u=\bigoplus_{a\in J\setminus\{0\}}
\mathfrak u[a]$ denote it and let $U$ be its simply connected real
group. Set $\mathfrak u_b=\bigoplus_{a\ge b}\mathfrak u[a]$.
This is a descending rational multidegree filtration with the
required bracket rule.

Sending generators to their basis vectors defines a rational
homomorphism $\rho:U\to G$. The killed words map to zero by the
filtration rule in $G$. It maps each positive layer onto $G_b$:
the generators tagged $b$ span $\mathfrak g_b$, and all weights
at least $b$ map into that layer. At zero it is onto because
$G$ is taut and generated by its coordinate groups $G_{(1,0)}$ and $G_{(0,1)}$.
Thus the polynomial lifts exist by published GTZ Corollary B.10,
applied to this layerwise-surjective filtered map, and can be
selected internally at the representative indices.

Choose a homogeneous rational basis of $\mathfrak u$ and let
$L_0$ be its integer span. There is a positive integer $D$ such
that $\Lambda=\exp(D L_0)$ is a lattice and $\rho(\Lambda)\subset
\Gamma$. Here are the scaling details. In this basis BCH is a
finite rational polynomial with linear part addition. For two
inputs in $D L_0$, terms of degree at least two lie in $D L_0$
once $D$ clears all the coefficient denominators. Inversion
negates the logarithm. Thus this set is a subgroup, discrete
in exponential coordinates. It is cocompact: order the basis
by increasing total weight and successively reduce each weight's
logarithm coordinates modulo $D$. Right multiplication by an
element of the corresponding homogeneous lattice leaves lower
weights unchanged, and at the current weight is ordinary
addition. The resulting bounded box of logarithms is a compact
set of representatives. The same argument applies to every
filtration layer and every homogeneous quotient. To impose
$\rho(\Lambda)\subset\Gamma$, use rational Mal'cev coordinates
of the second kind adapted to $\Gamma$, as in GTZ Section 3.
The coordinates of $\rho(\exp x)$ are finite rational
polynomials in the chosen coordinates of $x$, with zero constant
term. Taking $D$ divisible by their denominators makes the
coordinates of $\rho(\exp(D x))$ integral for $x\in L_0$.
This gives the required lattice inclusion. Increasing $D$ preserves
the BCH and compactness arguments.

For $i=0,1$, quotient $U$ by the connected subgroup with
Lie algebra the sum of homogeneous weights outside $J_i$.
These sums are ideals by the downset property. The image
lattices are the exponential lattices in the surviving
homogeneous bases. The pair of quotient maps separates points
of $U/\Lambda$: if two points have the same image in both,
the logarithm of their group difference lies in the projected
lattice in each quotient. Each homogeneous weight survives
in at least one quotient. Hence every coordinate of that
logarithm belongs to $D\mathbb Z$, so the difference lies in
$\Lambda$. This proves separation on the quotient spaces,
rather than just intersection-triviality of the two ideals.

The pullbacks of Lipschitz tests from the two compact quotients
generate a conjugation-closed algebra containing constants
and separating points. Stone--Weierstrass gives a uniform
approximation by finite sums of products of these tests.
The lifted tests of the shift family have a common Lipschitz
and sup bound on this fixed compact manifold, since the
map to $G/\Gamma$ is fixed and smooth. A standard finite
$\epsilon/4$ net for this bounded equicontinuous class gives
finitely many standard product lists, each chosen to tolerance
$\epsilon/4$. Select the first net entry within $\epsilon/4$
using an internal sup-norm comparison at representative indices.
Use a smaller net radius if needed. Thus the splitting error
is at most $\epsilon/2$, uniformly in $k,h,n$. Pad the lists
with zero factors to a standard common length.

In the $J_0$ quotient all surviving generators have zero first
weight and strictly positive second weight. The restricted
one-variable degree filtration is therefore taut: its degree-zero
and degree-one groups coincide. The highest layer is central by
the bracket law and the degree bound. The orbit is independent
of $h$, since the first-variable derivative lies in the trivial
first-coordinate group. Restrict to $h=0$ and apply published
Lemma E.5 with this taut one-variable geometry of degree $s-1$.
Use the ordinary functions constructed in its proof on printed p.~1353.
For each of the finitely many standard $J_0$ tests, fix ordinary Lipschitz
unit-vector tests with global vertical frequency and ordinary top-invariant
coefficient functions. Extend these functions to the ultrapower after making
these finite choices. The resulting tests have exact unit norm and
equivariance at all limit points and all limit top-group elements, and
internal selection from the finite lists preserves both identities.
Its test decomposition gives linear functionals of unit-vector
nilcharacters times degree-at-most-$(s-2)$ tests. Expand the
functionals in coordinates and absorb their coefficients and
lower-degree factors into the $J_1$ factors. The product still
has support $J_1$. Choose individual Fourier tolerances after
the finite lists and their sup bounds so that the added error
is at most $\epsilon/2$. On fixing $h,k$, the lower factor has
degree at most $s-2$ in $n$ by the $(0,1)$ derivative index.
All geometry lists are fixed and finite, and all orbit data
and choices are internal.
\end{proof}

\section{The two averages}
\begin{proposition}[LP03]
Under the hypotheses of GTZ Theorem 7.2, for standard $s\ge3$,
there are $\gg N^3$ additive quadruples in $H$ for which
\[
 \chi_{h_1}(n)\otimes\chi_{h_2}(n+h_1-h_4)
 \otimes\overline{\chi_{h_3}(n)}
 \otimes\overline{\chi_{h_4}(n+h_1-h_4)}
\]
is $(s-2)$-biased on $J$. The witnessing tests have standard
bounded complexity and positive standard correlation.
\end{proposition}
\begin{proof}
Extend $f$ by zero outside $J$. GTZ A.11 and A.12 select the
lower-degree correlation tests internally on the dense internal set
$H\subset I$. The relation includes a selected positive lower bound
$1/j$ with standard integer $j$. Select global representatives of the lower-degree tests and of the
restricted nilcharacters by the same internal selection, storing
the representations as part of the witness. Only their original
values on $J$ enter the unsmoothed identities. The bounded-complexity descriptions
are a countable union of internal sets. The boundedness principle
A.5 gives uniform standard complexity and a common positive lower
bound. Expand the bounded-dimensional tensor contraction in
coordinates. Retain a positive-density subset of $H$ with a fixed
witnessing coordinate pair by the standard finite pigeonhole.
Write $U_h,V_h$ for the conjugates of those coordinates of
$\chi(h,\cdot)$ and $\chi_h$. Then there are scalar lower-degree
tests $c_h$ such that
\[
 \left|\E_{n\in J}f(n+h)\overline{f(n)}U_h(n)V_h(n)c_h(n)\right|
 \ge\eta>0
\]
on the retained set, with standard uniform bounds. Align the phase
by multiplying $c_h$ by a unit scalar. Set $c_h=0$ outside that
set and extend $V_h$ boundedly there. This is an internal choice.
The real average over $h\in I$ is $\gg1$.

Apply Cauchy--Schwarz in $n$ to eliminate $\overline{f(n)}$.
It gives
\[
 \E_{n\in J}\left|\E_{h\in I}
       f(n+h)U_h(n)V_h(n)c_h(n)\right|^2\gg1.
\]
Expand the square and put $h'=h+k$, with
$K=[-2N,2N]\cap{}^*\Z$. The exact new parameter cutoff is
$1_{h+k\in I}$, degree zero in $n$. On $I\times K$ the new
normalized average is $|I|/|K|$ times the old one. With
$c_{h,k}=1_{h+k\in I}c_h\overline{c_{h+k}}$ this gives
\[
 \left|\E_{h\in I,k\in K,n\in J}
 f(n+h)\overline{f(n+h+k)}
 U_h(n)\overline{U_{h+k}(n)}
 V_h(n)\overline{V_{h+k}(n)}c_{h,k}(n)\right|\gg1.
\]
The $U$ product is a component of $D_k$. Take LP02 with error
less than half the correlation divided by a common bound for the
other factors. A term of the standard finite sum retains a positive
standard correlation. Absorb its lower factor into $c_{h,k}$.

For its surviving coordinate $\Psi_{k,a}(n)$ use the identity
\[
 \Psi_{k,a}(n)=\sum_b
 [\Psi_{k,a}(n)\overline{\Psi_{k,b}(n+h)}]\Psi_{k,b}(n+h).
\]
The identity uses the unit vector norm of $\Psi_k(n+h)$.
Individual coordinates need not have modulus one. The bracket
is a component of the conjugate transpose of the shift tensor.
The one-variable degree case of the project shift theorem gives
degree at most $s-2$, uniformly and internally in $h,k$, using the
finite geometries of LP02. Select a surviving term by the same
finite averaging argument. Absorb the bracket into $d_{h,k}$ and
write the remaining coordinate as $P_k(n+h)$. The lower test is
still exactly zero unless both $h,h+k$ are retained shifts.

Set $m=n+h$ and $M=[-2N,2N]\cap{}^*\Z$. The change to the
full $M$ average, with the exact cutoff, multiplies the
normalized correlation by $N/|M|$. We have
\[
 \left|\E_{m\in M,k\in K}
 f(m)\overline{f(m+k)}P_k(m)
 \E_{h\in I} 1_J(m-h)V_h(m-h)
 \overline{V_{h+k}(m-h)}d_{h,k}(m-h)\right|\gg1 .
\]
The outer factor depends only on $(m,k)$ and is bounded.
The second Cauchy--Schwarz in $(m,k)$ gives $\gg1$ for the
mean squared modulus of the inner average. Expand with independent
$h,h'\in I$. Put
\[
 (h_1,h_2,h_3,h_4)=(h,h'+k,h+k,h'),\qquad n=m-h.
\]
The shift substitution is injective, with inverse
$(h,h',k)=(h_1,h_4,h_3-h_1)$. It has $h_1+h_2=h_3+h_4$.
The first cutoff restricts $n$ to $J$, where $m=n+h\in M$
automatically. The second is $1_J(n+h_1-h_4)$. The four factors
in the new average are
\[
 V_{h_1}(n)V_{h_2}(n+h_1-h_4)
 \overline{V_{h_3}(n)}\overline{V_{h_4}(n+h_1-h_4)} .
\]
The lower factor is
$d_{h,k}(n)\overline{d_{h',k}(n+h_1-h_4)}$
times the remaining cutoff. The ratio $N/|M|$ is comparable
to one. The parameter vanishing is exact unless all four shifts
are retained. This proof keeps zero extension throughout.

To absorb the cutoff, choose fixed $0<\lambda<1/4$.
On the circle $\mathbb R/16\Z$ take a piecewise linear test
equal to one on $[\lambda,1-\lambda]$, zero off
$[-\lambda,1+\lambda]$ in $[-8,8)$, between zero and one,
with Lipschitz bound $O(1/\lambda)$. Evaluate it at
$(n+h_1-h_4)/N$. This is a one-step nilsequence in $n$,
with fixed geometry and internal orbit. Its argument lies in
$[-2,3]$ and does not wrap. The normalized mean absolute error
against the cutoff is $O(\lambda+1/N)$ uniformly in all shifts:
only two endpoint bands of total $O(\lambda N+1)$ integers
differ. Multiplying by bounded factors changes only the constant.
Choose $\lambda$ after the positive correlation bound so the
error loses at most half. Since $s\ge3$, the one-step test is
absorbed into degree at most $s-2$. Parameter vanishing persists.

If the outer correlation has modulus at least $\beta>0$ and
each inner average has modulus at most $B$, at least a fraction
$\beta/(2B)$ of the inner averages have modulus at least
$\beta/2$. The shift triple box has size comparable to $N^3$.
By injectivity, these yield $\gg N^3$ additive quadruples in $H$.
Conjugating the scalar correlations undoes the initial choice of
$V_h$ and preserves lower degree. The resulting scalar correlations
are coordinates of the claimed tensor, giving bias on the
original interval $J$.
\end{proof}

\section{Scope}
The published proof inputs used here are GTZ Appendix A, Corollary
B.10, E.5, the correlation conventions, and the stated
hypotheses of Theorem 7.2. The project shift input and these exact
applications await separate nonauthor review. This unit supplies the
additive-quadruple result, not the subsequent frequency analysis,
the other affected symbol rules, or the final GI theorem.
\end{document}
